% !TeX program = xelatex
%%%%%%%%%%%%%%%%%%%%%%%%%%%%%%%%%%%%%%%%%
% Twenty Seconds Resume/CV
% LaTeX Template
% Version 1.0 (14/7/16)
%
% Original author:
% Carmine Spagnuolo (cspagnuolo@unisa.it) with major modifications by
% Vel (vel@LaTeXTemplates.com) and Harsh (harsh.gadgil@gmail.com)
%
% License:
% The MIT License (see included LICENSE file)
%
%%%%%%%%%%%%%%%%%%%%%%%%%%%%%%%%%%%%%%%%%

%----------------------------------------------------------------------------------------
%	PACKAGES AND OTHER DOCUMENT CONFIGURATIONS
%----------------------------------------------------------------------------------------

\documentclass[letterpaper]{twentysecondcv} % a4paper for A4

% Command for printing skill overview bubbles
\newcommand\skills{
    \begin{itemize}
        \item Neovim
        \item Linux
    \end{itemize}
}

\education{
\begin{twenty}
    \educationitem{2017 - 2021}{\textbf{MSc. in Complex Adaptive Systems} \\ Chalmers University of Technology \\ \textit{Gothenburg, Sweden}}\\
    \educationitem{2013 - 2017}{\textbf{Information Technology and Software Engineering} \\ Chalmers University of Technology \\ \textit{Gothenburg, Sweden}}
\end{twenty}
}

% Programming skill bars
%\programming{{C $\textbullet$ C++  $\textbullet$ R / 3}, {JavaScript $\textbullet$ SQL $\textbullet$ \large \LaTeX / 3.5}, {Scala $\textbullet$ Java $\textbullet$ Python / 5}}
\renewcommand{\programming}{
    \textbf{Languages}
    \begin{itemize}
        \item C++
        \item Java
        \item Python
        \item Rust
        \item Bash/Shell scripting
    \end{itemize}
    \textbf{Tools}
    \begin{itemize}
        \item Editors: Vim/Neovim, VSCode
        \item Git
        \item Platforms: Linux, Unix, Docker, LXC
    \end{itemize}
}

% Languages
\languages{
    \languageitem{Chinese (Mandarin)}{Native}
    \hline
    \languageitem{Swedish}{Native}
    \hline
    \languageitem{English}{Fluent}
}

% Projects text

%----------------------------------------------------------------------------------------
%	 PERSONAL INFORMATION
%----------------------------------------------------------------------------------------
% If you don't need one or more of the below, just remove the content leaving the command, e.g. \cvnumberphone{}

\cvname{Henry Yang} % Your name
\cvjobtitle{IT Engineer} % Job
% title/career
\profilepic{img/trophy.png} %placeholder
\cvlinkedin{/in/hyanggbg}
\cvgithub{Xiaoming94}
\cvnumberphone{(+46) 0704 50 84 51} % Phone number
\cvmail{mail@zyming.me} % Email address

%----------------------------------------------------------------------------------------

\begin{document}

\makeprofile % Print the sidebar
\small

%----------------------------------------------------------------------------------------
%    ABOUT ME
%----------------------------------------------------------------------------------------
\section{About me}
Henry is a developer and engineer with experience in Telecom.
Henry commands a wide array of skills from programming and software engineering, to managing server infrastructure, and solution systemisation.
He is dedicated to his work and always puts in the necessary effort to solve a problem.
He is also a fast learner and doesn't mind spending time on research and studying to find a solution on extra complex stuff.

With a degree in Master of Science in Complex Adaptive Systems from Chalmers University of Technology, Henry is leveraging his knowledge from CAS in the problem solving process.
During his days in the University, Henry has picked up various technical skills from extra curricular activities which he are still making use of in the daily.
He is always up for sharing his knowledge with his fellow colleagues to help them solving the problem they are facing.

%---------------------------------------------------------------------------------------
%	 EXPERIENCE
%----------------------------------------------------------------------------------------

\section{Experience}

\begin{twenty} % Environment for a list with descriptions
\twentyitem
        {2021 -}
        {Ongoing}
        {Consultant}
        {\href{https://nexergroup.com/sv/}{Nexer Engineering 1 AB}}
        {\scriptsize Employee}
        {\small
            Consultant company that is part of the larger Nexer Group AB.
            Henry got his assignment at Ericsson from here.

            Henry also engages in extra curricular activities at Nexer.
            Notably, he created study material for the programming language rust.
            He also held introductory presentation and workshop for to help people get started with this programming language.
        }\\

\twentyitem
        {2021 -}
        {Ongoing}
        {SW Developer}
        {\href{https://ericsson.com}{Ericsson AB}}
        {\scriptsize Assignment}
        {\small
            At ericsson, Henry works as a C++ feature developer for the Ericsson Radio Base Stations.
            During his time at Ericsson, Henry and his team have been involved in various key features in Ericssons cellular system.

            Most notably, he took part in developing the fault management system of NG-OLLS, Ericssons low-level ORAN solution.
            Apart from taking part in the development process, Henry contributed to the research and high-level design of the fault handling process.
            \\
            \skilltag{C++} \skilltag{linux} \skilltag{googletest} \skilltag{java} \skilltag{bash} \skilltag{moshell} \skilltag{CI/CD}
        \skilltag{Test Driven Development} \skilltag{OOP}
        }

	%\twentyitem{<dates>}{<title>}{<location>}{<description>}
\end{twenty}

%----------------------------------------------------------------------------------------
%    EDUCATION
%----------------------------------------------------------------------------------------
\section{Education}

\education

%----------------------------------------------------------------------------------------
%	 RESEARCH
%----------------------------------------------------------------------------------------

\section{Publications}
L. Antonie, H. Gadgil, G. Grewal, and K. Inwood, “Historical Data Integration - A Study of WWI Canadian Soldiers,” in 2016 IEEE 16th International Conference on Data Mining Workshops (ICDMW), pp. 186-193, IEEE, 2016. \vspace{2mm}

\end{document}
